% FILE: 03_architecture_and_primitives.tex
% Project: experiments-visualizer-nextjs
% Role: web-based-pi-visualization-surface

\section{Architecture and Primitives}
\label{sec:experiments-visualizer-nextjs-arch}

\subsection{Core Objects}
\label{sec:experiments-visualizer-nextjs-arch-objects}

The visualizer-nextjs project defines its core objects through the CSS class vocabulary in
\texttt{globals.css}. These classes are not merely styling tokens; each one names a
structural object in the process-intelligence domain that the dashboard is responsible for
rendering. The five principal object families are:

\begin{description}
  \item[\textbf{Petri Net Elements}]
    \texttt{.petri-place}, \texttt{.petri-trans}, \texttt{.petri-arc},
    \texttt{.petri-place.active}, \texttt{.petri-trans.enabled}, \texttt{.place-token-text}.
    These classes define the visual atoms of a Petri net SVG panel: places (circles),
    transitions (rectangles), arcs (directed edges), active places (holding tokens),
    enabled transitions (fireable given current marking), and the token count label inside
    an active place. Together they encode the full marking-visible state of a Petri net
    replay step.

  \item[\textbf{A{*} Alignment Blocks}]
    \texttt{.alignment-block.sync}, \texttt{.alignment-block.log},
    \texttt{.alignment-block.model}. These three classes correspond precisely to the three
    move types in A{*} process alignment: a synchronous move (event matches model transition),
    a log move (event has no counterpart in the model), and a model move (model transition
    has no counterpart in the event log). The alignment block sequence is the primary
    per-trace conformance artefact.

  \item[\textbf{EWMA Drift Detection}]
    \texttt{.ewma-canvas}, \texttt{.alarm-banner}, \texttt{.alarm-banner.fired}. These
    classes define the drift-monitoring panel: a canvas element for the EWMA time series
    and an alarm banner that transitions to the \texttt{fired} state when the EWMA statistic
    crosses the control limit. The binary fired/unfired state corresponds to the drift alarm
    signal produced by the upstream EWMA computation.

  \item[\textbf{Cryptographic Ledger}]
    \texttt{.ledger-block-card.healthy}, \texttt{.ledger-block-card.tampered},
    \texttt{.ledger-chain-arrow.broken}, \texttt{.receipt-json-body},
    \texttt{.verification-console}. These classes define the receipt ledger explorer.
    A block card is healthy when its hash-chain link is verified and tampered when
    verification fails; a chain arrow is broken when the link between two consecutive
    blocks cannot be verified. The receipt JSON body and verification console display
    the raw receipt blob and the pass/fail verdict.

  \item[\textbf{Dashboard Layout}]
    \texttt{.dashboard-grid} (320px sidebar + 1fr main area), \texttt{.panel},
    \texttt{.global-metrics-bar}. These classes define the top-level layout container
    that positions the four evidence panels in a single-page dashboard.
\end{description}

\subsection{Admissible Transitions}
\label{sec:experiments-visualizer-nextjs-arch-transitions}

AUTHOR THESIS: The CSS class algebra implies a set of admissible UI state transitions that
correspond to lawful transitions in the underlying process-intelligence state machine.

\begin{itemize}
  \item A \texttt{.petri-place} may transition to \texttt{.petri-place.active} only when a
        token is deposited by a fired transition; it may return to inactive only when a token
        is consumed.
  \item A \texttt{.petri-trans} may bear the \texttt{.enabled} modifier only when all input
        places hold sufficient tokens; it loses \texttt{.enabled} when any input place is
        drained.
  \item An \texttt{.alignment-block} must carry exactly one of the three move-type modifiers
        (\texttt{sync}, \texttt{log}, \texttt{model}); no block may carry zero or more than one.
  \item An \texttt{.alarm-banner} transitions to \texttt{.fired} only when the EWMA statistic
        exceeds the control limit; it may not oscillate without a corresponding change in the
        underlying time series value.
  \item A \texttt{.ledger-block-card} may transition from \texttt{healthy} to \texttt{tampered}
        when hash verification fails; the reverse transition is not admissible without a
        re-verification event being recorded in the receipt chain.
\end{itemize}

\subsection{Boundaries}
\label{sec:experiments-visualizer-nextjs-arch-boundaries}

The visualizer-nextjs project is bounded by the following architectural constraints established
by the project's source files and the broader research program:

\begin{itemize}
  \item \textbf{Framework:} Next.js 16.2.6 App Router with Turbopack. The \texttt{AGENTS.md}
        file explicitly warns that this version contains breaking changes relative to prior
        training data and instructs any authoring agent to read
        \texttt{node\_modules/next/dist/docs/} before writing code.
  \item \textbf{Language:} TypeScript 5 with strict mode (\texttt{tsconfig} confirmed).
  \item \textbf{Styling:} Tailwind CSS v4 with \texttt{@tailwindcss/postcss}; CSS custom
        properties for the color system; no CSS-in-JS library.
  \item \textbf{Font system:} Geist and Geist Mono via the Vercel font system, loaded in
        \texttt{src/app/layout.tsx}.
  \item \textbf{Scope:} The visualizer is a read-only observation surface. It does not write
        to the event log, modify process models, or issue receipts. Those responsibilities
        belong to upstream Rust and Python components.
\end{itemize}

\subsection{Failure Modes Refused}
\label{sec:experiments-visualizer-nextjs-arch-failures}

The architecture of the CSS design system explicitly refuses several failure modes by making
them unrepresentable:

\begin{itemize}
  \item \textbf{Ambiguous alignment moves:} The three-modifier system for alignment blocks
        (\texttt{sync}, \texttt{log}, \texttt{model}) makes it structurally impossible to
        render an alignment block with no move type or with a fabricated move type not
        recognized by the A{*} algorithm.
  \item \textbf{Silent tamper:} The \texttt{.ledger-block-card.tampered} class and
        \texttt{.ledger-chain-arrow.broken} class make hash-chain failure visually
        unmissable; the design system does not provide a class for a tampered block that
        appears healthy.
  \item \textbf{Fired alarm without visual indicator:} The \texttt{.alarm-banner.fired}
        modifier has no silent state --- a fired alarm is always rendered as a distinct
        visual element.
  \item \textbf{Unimplemented panel rendered as empty:} AUTHOR THESIS: The correct failure
        mode for a panel whose data source is not connected is an explicit error state, not
        an empty render. The CSS design system does not yet define an error-state class
        for missing data; this is an open gap.
\end{itemize}
